\documentclass[12pt,a4paper]{article}
% ============================================================
%  Оформление отчёта. Обычно здесь ничего менять не нужно.
%  Собирается pdfLaTeX (по умолчанию, в т. ч. в Overleaf),
%  XeLaTeX или LuaLaTeX.
% ============================================================

\usepackage{iftex}
\ifPDFTeX
  \usepackage[T2A]{fontenc}
  \usepackage[utf8]{inputenc}
  \usepackage{tempora}                 % Times-подобный шрифт с кириллицей
\else
  \usepackage{fontspec}
  \IfFontExistsTF{Times New Roman}
    {\setmainfont{Times New Roman}}
    {\IfFontExistsTF{CMU Serif}
      {\setmainfont{CMU Serif}}
      {\IfFontExistsTF{Liberation Serif}
        {\setmainfont{Liberation Serif}}
        {\setmainfont{DejaVu Serif}[AutoFakeSlant=0.2]}}}
  \IfFontExistsTF{Consolas}
    {\setmonofont{Consolas}[Scale=MatchLowercase]}
    {\IfFontExistsTF{Courier New}
      {\setmonofont{Courier New}}
      {\IfFontExistsTF{CMU Typewriter Text}
        {\setmonofont{CMU Typewriter Text}}
        {\setmonofont{DejaVu Sans Mono}[Scale=MatchLowercase]}}}
\fi
\IfFileExists{russianb.ldf}
  {\usepackage[english,russian]{babel}}
  {\usepackage[english]{babel}\babelprovide[import,main]{russian}}

% --- Поля и абзацы (как в исходном doc-шаблоне курса) ---
\usepackage[a4paper, margin=2.54cm]{geometry}
\setlength{\parindent}{0pt}
\setlength{\parskip}{4pt plus 1pt}
\frenchspacing
% Запас для переноса текста при сборке без русских шаблонов переносов.
\setlength{\emergencystretch}{0.5\linewidth}

% --- Заголовки разделов без номеров ---
\usepackage{titlesec}
\setcounter{secnumdepth}{0}
\titleformat{\section}{\large\bfseries}{}{0pt}{}
\titleformat{\subsection}{\normalsize\bfseries}{}{0pt}{}
\titlespacing*{\section}{0pt}{14pt plus 4pt minus 2pt}{4pt plus 2pt}
\titlespacing*{\subsection}{0pt}{10pt plus 2pt minus 2pt}{3pt plus 1pt}

% --- Таблицы, рисунки, подписи ---
\usepackage[table]{xcolor}
\usepackage{graphicx}
\usepackage{tabularx}
\usepackage{float}
\usepackage{caption}
\DeclareCaptionLabelSeparator{spacedendash}{~-- }
\captionsetup{labelsep=spacedendash, font=small, singlelinecheck=false, justification=raggedright}
\captionsetup[figure]{justification=centering, singlelinecheck=true}
\captionsetup[table]{position=top, skip=4pt}
\DeclareCaptionType[fileext=lol]{listing}[Листинг][Листинги]
\captionsetup[listing]{position=top, skip=4pt, hypcap=false}
\addto\captionsrussian{%
  \renewcommand{\figurename}{Рисунок}%
  \renewcommand{\tablename}{Таблица}%
}
\newcommand{\thead}[1]{\multicolumn{1}{|c|}{\cellcolor{gray!20}\textbf{#1}}}
% Колонки таблиц без растягивания пробелов: L – резиновая, P{ширина} – фиксированная
\newcolumntype{L}{>{\raggedright\arraybackslash}X}
\newcolumntype{P}[1]{>{\raggedright\arraybackslash}p{#1}}

% --- Код и вывод терминала ---
% code     – фрагмент исходного кода;
% terminal – вывод программы, strace и т. п.;
% \inputterminal{файл} / \inputcode{файл} – то же, но из файла.
\usepackage{needspace}
\usepackage{fvextra}
\fvset{breaklines=true, breakanywhere=true,
       frame=leftline, framerule=1.5pt, framesep=6pt,
       rulecolor=\color{gray!60}, xleftmargin=4pt}
\DefineVerbatimEnvironment{code}{Verbatim}{fontsize=\small}
\DefineVerbatimEnvironment{terminal}{Verbatim}{fontsize=\footnotesize}
\newcommand{\inputcode}[2][]{\VerbatimInput[fontsize=\small,#1]{#2}}
\newcommand{\inputterminal}[2][]{\VerbatimInput[fontsize=\footnotesize,#1]{#2}}
% Листинг с подписью сверху:
%   \begin{listingblock}{Подпись} \begin{code}...\end{code} \end{listingblock}
\newenvironment{listingblock}[1]
  {\par\medskip\Needspace{5\baselineskip}\captionof{listing}{#1}\nopagebreak}
  {\par}

% Приложение: с новой страницы, заголовок «Приложение А. ...»,
% листинги внутри нумеруются А.1, А.2, ...
\newcounter{appendixno}
\newcommand{\appendixsection}[2]{%
  \clearpage
  \section{Приложение #1. #2}%
  \stepcounter{appendixno}%
  \setcounter{listing}{0}%
  \renewcommand{\thelisting}{#1.\arabic{listing}}%
  \renewcommand{\theHlisting}{app\arabic{appendixno}.\arabic{listing}}}

% --- Ссылки ---
\usepackage[hidelinks]{hyperref}
\usepackage{xurl}

% --- Подсказки шаблона: серый курсив, удалить перед сдачей ---
\newcommand{\hint}[1]{\leavevmode{\color{gray}\itshape #1}}
% Рамка-заглушка вместо рисунка
\newcommand{\placeholder}[2][4.5cm]{%
  \fbox{\parbox[c][#1][c]{0.95\linewidth}{\centering\hint{#2}}}}


\newcommand{\ReportLab}{3}
\newcommand{\ReportStudent}{Федосеев Николай Владимирович}
\newcommand{\ReportGroup}{М8О-204БВ-25}
\newcommand{\ReportTeacher}{Н. А. Гапонов}
\newcommand{\ReportDate}{22 октября 2026 г.}
\newcommand{\ReportYear}{2026}
\newcommand{\ReportCourse}{Операционные системы и архитектура компьютера}

\begin{document}
% Титульный лист. Данные берутся из команд \Report... в report.tex.
\begin{titlepage}
  \centering\large

  МОСКОВСКИЙ АВИАЦИОННЫЙ ИНСТИТУТ\\
  (НАЦИОНАЛЬНЫЙ ИССЛЕДОВАТЕЛЬСКИЙ УНИВЕРСИТЕТ)\par
  \vspace{18pt}
  Институт №8 «Компьютерные науки и прикладная математика»\\
  Кафедра 806 «Вычислительная математика и программирование»\par

  \vspace{\stretch{1.2}}

  {\bfseries Лабораторная работа №\,3 по курсу\\
    «Операционные системы и архитектура компьютера»\par}

  \vspace{\stretch{1.6}}

  \begin{flushright}
    \begin{tabular}{r@{\hspace{0.6em}}l}
      Выполнил:      & Федосеев Николай Владимирович \\[2pt]
      Группа:        & М8О-204БВ-25         \\[2pt]
      Преподаватель: & Н. А. Гапонов   \\[2pt]
      Дата:          & 22.10.2026          \\[2pt]
      Оценка:        & \rule[-2pt]{4.2cm}{0.4pt} \\
    \end{tabular}
  \end{flushright}

  \vspace{\stretch{1}}

  Москва, 2026
\end{titlepage}
\setcounter{page}{2}


\section{Условие}
\textbf{Цель работы:} получить практические навыки управления потоками в ОС и синхронизации между ними.

\textbf{Задание, вариант 9.}

Рассчитать детерминант матрицы, используя определение детерминанта.

Программа должна быть написана на языке Си и использовать стандартные средства создания потоков ОС. Максимальное количество одновременно работающих потоков задаётся ключом запуска. Необходимо показать количество потоков средствами ОС, исследовать ускорение и эффективность алгоритма при изменении входных данных и числа потоков, объяснить результаты.

\section{Метод решения}
Определитель вычисляется по определению: перебираются перестановки столбцов, для каждой берётся по одному элементу из каждой строки и столбца. Их произведение прибавляется или вычитается в зависимости от чётности перестановки:
\[
\det A = \sum_{\sigma \in S_n} (-1)^{\mathrm{inv}(\sigma)} \prod_{i=0}^{n-1} a_{i,\sigma(i)}.
\]
Здесь $\mathrm{inv}(\sigma)$ - число инверсий: пар, в которых левый элемент перестановки больше правого. Всего имеется $n!$ слагаемых. Работа прямого перебора имеет порядок $O(n\cdot n!)$; потоки распределяют эту работу, но не устраняют факториальный рост.

При $n\geq2$ вычисление разбивается на задачи (\texttt{job} в коде). Каждая задача перебирает перестановки с фиксированными первыми двумя элементами. Всего таких задач $n(n-1)$, каждая перебирает $(n-2)!$ вариантов хвоста. При $n=1$ задача одна. Функция \texttt{prepare\_permutation} заполняет выбранное начало и отсортированный хвост, затем определяет начальный знак по инверсиям. \texttt{permutation\_next} переставляет хвост и переключает знак при каждом обмене.

Исполнитель с номером $i$ получает задания $i$, $i+p$, $i+2p$ и так далее, где $p$ - фактическое число исполнителей. Главный поток тоже считает: создаются $p-1$ дополнительных потоков. Число исполнителей ограничивается числом заданий, но не количеством ядер.

Матрица общая и во время расчёта не меняется. У каждого исполнителя свой массив перестановки. Сумма задания записывается в отдельную ячейку \texttt{sums[job]}, которую другие исполнители не изменяют. После завершения потоков главный складывает ячейки в порядке номеров заданий. Общая очередь и мьютекс для вычислений не нужны: задачи распределены заранее, а ожидание и получение результатов выполняются через \texttt{pthread\_join}.

\section{Описание программы}
\subsection{Структура файлов}
\begin{table}[H]
\caption{Основные файлы проекта}
\begin{tabularx}{\linewidth}{|P{4.6cm}|L|}
\hline
\thead{Файл} & \thead{Назначение} \\ \hline
\texttt{src/main.c} & Аргументы, открытие и закрытие файла, замер времени и вывод. \\ \hline
\texttt{src/matrix.h/.c} & Структура матрицы, статусы, выделение и освобождение памяти. \\ \hline
\texttt{src/matrix\_txt.h/.c} & Чтение матрицы из текстового потока. \\ \hline
\texttt{src/permutation.h/.c} & Следующая перестановка и обновление знака. \\ \hline
\texttt{src/determinant.h/.c} & Разбиение на задания, рабочая функция, запуск и ожидание потоков, сбор результата. \\ \hline
\texttt{src/generate.py} & Простой генератор TXT: размер и $n^2$ случайных чисел. \\ \hline
\texttt{src/CMakeLists.txt} & Сборка программы C17 с POSIX threads. \\ \hline
\end{tabularx}
\end{table}

\subsection{Основные типы данных и функции}
\texttt{struct Matrix} хранит размер \texttt{size} и указатель \texttt{data} на непрерывный массив \texttt{int64\_t}. Элемент строки \texttt{row} и столбца \texttt{column} имеет индекс \texttt{row * size + column}. Размеры и индексы имеют тип \texttt{size\_t}.

\texttt{matrix\_init} выделяет память пустой матрице, \texttt{matrix\_destroy} освобождает её и обнуляет поля. \texttt{matrix\_read\_txt} читает размер и элементы во временную матрицу, а наружу передаёт её только после успешного чтения. Открытый поток принадлежит \texttt{main}.

\texttt{struct Worker} содержит адрес матрицы, номер исполнителя, количество исполнителей и заданий, общий массив сумм и признак успеха. \texttt{worker\_run(void *argument)} получает адрес этой структуры, выделяет собственную перестановку и выполняет свои задания.

\texttt{determinant\_compute} выделяет общие рабочие массивы, подготавливает контексты, создаёт дополнительные потоки, выполняет работу исполнителя 0 и присоединяет созданные потоки. Возвращает \texttt{true} или \texttt{false}. \texttt{struct Calculation} получает определитель и фактическое число исполнителей только при успехе.

\subsection{Потоки и системные вызовы}
\begin{table}[H]
\caption{Основные операции программы}
\begin{tabularx}{\linewidth}{|P{4.6cm}|L|}
\hline
\thead{Функция в коде} & \thead{Назначение} \\ \hline
\texttt{pthread\_create} & Запускает \texttt{worker\_run} в новом потоке процесса. \\ \hline
\texttt{pthread\_join} & Ожидает завершения конкретного рабочего потока. \\ \hline
\texttt{clock\_gettime} & Получает время \texttt{CLOCK\_MONOTONIC} до и после расчёта. \\ \hline
\texttt{fopen}, \texttt{fscanf}, \texttt{fclose} & Открывают, читают и закрывают текстовый файл. \\ \hline
\texttt{malloc}, \texttt{free} & Выделяют и освобождают матрицу и рабочие массивы. \\ \hline
\texttt{printf}, \texttt{fflush} & Выводят результат и проверяют сброс буфера вывода. \\ \hline
\end{tabularx}
\end{table}

\subsection{Ключевые фрагменты кода}
Каждый исполнитель получает непересекающиеся номера заданий.
\begin{listingblock}{Распределение заданий (determinant.c)}
\inputcode[firstline=60,lastline=62]{../src/determinant.c}
\end{listingblock}

Перед запуском потоков все контексты уже заполнены. При успешном \texttt{pthread\_create} счётчик \texttt{created} увеличивается. Затем главный поток сам вызывает \texttt{worker\_run} для контекста 0. При ошибке создания этот вызов пропускается, но уже запущенные потоки всё равно присоединяются.
\begin{listingblock}{Запуск дополнительных потоков (determinant.c)}
\inputcode[firstline=130,lastline=146]{../src/determinant.c}
\end{listingblock}

В произведениях и суммах используется \texttt{int64\_t}. Каждое умножение, сложение и вычитание проверяется через \texttt{\_\_builtin\_*\_overflow}.

\subsection{Обработка ошибок}
Проверяются аргументы запуска, открытие, чтение и закрытие файла, размеры выделений, результат \texttt{malloc}, ошибки потоков, арифметика, получение времени и вывод. Размеры выделений проверяются до умножения. При частичном отказе \texttt{malloc} успешно выделенные массивы освобождаются в общем месте очистки.

При ошибке \texttt{pthread\_create} программа дожидается только действительно созданных потоков и возвращает неуспех. При арифметической ошибке исполнитель завершает свою работу, остальные заканчивают свои задания; итоговый определитель не публикуется. При ошибке \texttt{pthread\_join} процесс завершается через \texttt{\_Exit}: освобождать память, которой ещё может пользоваться поток, небезопасно.

Загрузчик рассчитан на подготовленный числовой ввод: значения должны помещаться в \texttt{int64\_t}. \texttt{fscanf} здесь не обеспечивает безопасную проверку произвольно больших числовых записей. Переполнение промежуточного произведения или суммы считается ошибкой, даже если после математических сокращений итог мог бы поместиться.

\section{Сборка и запуск}
Используется CMake не ниже 4.3, стандарт C17. Макрос \texttt{\_POSIX\_C\_SOURCE=200809L} открывает объявления POSIX-функций. Потоки подключаются через \texttt{Threads::Threads}.
\begin{terminal}
cmake -S src -B build -DCMAKE_BUILD_TYPE=Debug
cmake --build build --parallel
./build/determinant --threads 3 src/examples/matrix.txt
\end{terminal}

Файл содержит положительный размер, затем ровно $n^2$ элементов. Для матрицы ниже определитель равен $-3$.
\begin{terminal}
3
1 2 3
4 5 6
7 8 10
\end{terminal}

Флаг \texttt{\-\-threads} задаёт предел исполнителей вместе с главным потоком. Программа выводит определитель, фактическое число исполнителей и время. Замер \texttt{CLOCK\_MONOTONIC} охватывает весь \texttt{determinant\_compute}: выделение рабочей памяти, создание потоков, вычисления, ожидание, сложение сумм и очистку. Чтение файла и вывод исключены из времени подсчёта.

\section{Демонстрация работы strace}
Программа запущена из корня проекта командой:
\begin{terminal}
strace -f -yy -s 256 -o report/inc/strace-3x3.txt ./build/determinant --threads 3 src/examples/matrix.txt
\end{terminal}

\texttt{-f} включает рабочие потоки, \texttt{-yy} подписывает дескрипторы, \texttt{-s 256} увеличивает лимит отображаемых строковых данных, \texttt{-o} сохраняет трассу в файл. Полная трейс приведена в приложении Б. Номера строк ниже относятся к этому файлу.

\subsection{Чтение матрицы}
Загрузка библиотек занимает начало трассы. В строке 32 \texttt{openat} открывает \texttt{src/examples/matrix.txt} и возвращает дескриптор 3. На строке 34 \texttt{read} получает все 21 байт файла: размер и девять элементов. На строке 35 чтение возвращает 0, то есть EOF. На строке 36 дескриптор закрывается. Несколько вызовов \texttt{fscanf} используют буфер libc, поэтому отдельного \texttt{read} для каждого числа нет.

\subsection{Создание потоков}
Главный поток имеет TID 508983. В строках 42 и 47 два успешных \texttt{clone3} создают потоки с TID 508984 и 508985. Вместе с главным это три исполнителя, как задано ключом запуска.
\begin{listingblock}{Первый рабочий поток: системный вызов clone3}
\inputterminal[firstline=42,lastline=42]{inc/strace-3x3.txt}
\end{listingblock}

Флаг \texttt{CLONE\_VM} означает общую память, \texttt{CLONE\_FILES} - общую таблицу дескрипторов, \texttt{CLONE\_THREAD} - принадлежность одному процессу. Поэтому это потоки одного процесса, а не три отдельных процесса. \texttt{CLONE\_SETTLS} задаёт локальное хранилище потока. Перед созданием каждого рабочего потока библиотека выделяет память под его стек через \texttt{mmap} с флагом \texttt{MAP\_STACK}. Затем вызовом \texttt{madvise} с параметром \texttt{MADV\_GUARD\_INSTALL} первые 4096 байт выделенной области отмечаются как защитные: обращение к ним вызывает \texttt{SIGSEGV}, что помогает обнаружить выход за границу стека. Эти действия видны в строках 39--40 и 44--45.

\subsection{Ожидание и завершение}
На строке 49 главный поток блокируется в \texttt{futex}; адрес совпадает с \texttt{child\_tid} первого \texttt{clone3}. После завершения рабочих потоков вызов возвращает 0 на строке 68.
\begin{listingblock}{Начало ожидания и возврат из futex}
\inputterminal[firstline=49,lastline=49]{inc/strace-3x3.txt}
\inputterminal[firstline=68,lastline=68]{inc/strace-3x3.txt}
\end{listingblock}

В коде нет мьютекса: \texttt{futex} используется библиотекой потоков при ожидании завершения. Второй поток уже завершился к проверке его \texttt{join}, поэтому отдельного ожидания в трейсе не видно. Пометки \texttt{unfinished} и \texttt{resumed} - части записи одного вызова, между которыми в лог попали события других потоков.

Рабочие потоки завершаются через \texttt{exit(0)} на строках 61 и 67. Затем главный выводит результат:
\begin{terminal}
Determinant: -3
Threads: 3
Time: 0.003735860 s
\end{terminal}

Это видно по трём \texttt{write} на строках 72--74. На строке 75 \texttt{exit\_group(0)} завершает процесс.

\section{Исследование ускорения и эффективности}
\subsection{Условия измерений}
Использована одна матрица $13\times13$ из \texttt{src/examples/matrix-13.txt}. Во всех девяти запусках получен одинаковый определитель 7736. Изменялось только число исполнителей: 1, 2, 3, 4, 5, 6, 7, 8 и 12. Для каждого значения передан один замер: среднее и медиана не вычислялись. Исходный вывод сохранён в приложении В.

Процессор - AMD Ryzen 5 7535HS: 6 физических ядер, 12 логических CPU, по два аппаратных потока на ядро (SMT). Система - Fedora Linux 44, архитектура x86-64. Использован каталог \texttt{build}; его конфигурация на момент подготовки отчёта - \texttt{Debug}, флаги \texttt{-g} без \texttt{-O3}. Запуски бенчмарков выполнены без \texttt{strace}.

Для $n=13$ имеется $13!=6\,227\,020\,800$ слагаемых, распределённых по $13\cdot12=156$ заданиям. В каждом задании $11!=39\,916\,800$ перестановок. По сравнению с шестью слагаемыми матрицы $3\times3$ это объясняет большой рост объёма работы; отдельные замеры разных размеров без трассировки не предоставлены.

\subsection{Результаты}
Ускорение и эффективность рассчитаны относительно однопоточного запуска:
\[
S_p=\frac{T_1}{T_p},\qquad E_p=\frac{S_p}{p}\cdot100\%,\qquad T_1=244.303998558\ \mathrm{s}.
\]
Здесь $p$ включает главный поток. Эффективность показывает отношение ускорения к числу исполнителей.

\begin{table}[H]
\caption{Время, ускорение и эффективность для матрицы $13\times13$}
\begin{tabularx}{\linewidth}{|P{1.6cm}|L|L|L|}
\hline
\thead{Потоков} & \thead{Время, с} & \thead{Ускорение} & \thead{Эффективность, \%} \\ \hline
1 & 244.303998558 & 1.000 & 100.00 \\ \hline
2 & 123.723745589 & 1.975 & 98.73 \\ \hline
3 & 85.636897709 & 2.853 & 95.09 \\ \hline
4 & 62.173046865 & 3.929 & 98.24 \\ \hline
5 & 52.815789465 & 4.626 & 92.51 \\ \hline
6 & 43.158549481 & 5.661 & 94.34 \\ \hline
7 & 53.069405205 & 4.603 & 65.76 \\ \hline
8 & 45.312672544 & 5.392 & 67.39 \\ \hline
12 & 41.237441684 & 5.924 & 49.37 \\ \hline
\end{tabularx}
\end{table}

\subsection{Разбор результатов}
При переходе от одного к четырём исполнителям время уменьшилось с 244.304 до 62.173 с, ускорение составило 3.929. Для шести исполнителей получено 43.159 с и ускорение 5.661. В этих запусках рост числа исполнителей близок к использованию отдельных физических ядер.

Минимальное время среди переданных замеров - 41.237 с при 12 исполнителях, ускорение 5.924. По сравнению с шестью исполнителями время уменьшилось только примерно на 4.45\%. Двенадцать логических CPU не равны двенадцати независимым ядрам: два аппаратных потока одного ядра разделяют его ресурсы. Поэтому эффективность при 12 исполнителях составляет 49.37\%, хотя ускорение относительно одного потока почти шестикратное.

Зависимость не монотонна: семь исполнителей дали 53.069 с, восемь - 45.313 с, то есть оба запуска медленнее шестипоточного. Возможные причины - размещение потоков по физическим ядрам и SMT, фоновые задачи, частоты и нагрев процессора. Также задания неделимы: 156 делится на 6 и 12 без остатка, а при 7 и 8 исполнителях число заданий на исполнителя различается.

Числа 7, 8 и 12 превышают количество физических ядер, но не число логических CPU моего ноутбука. Превышение 12 логических CPU в этих измерениях не исследовалось. При дальнейшем увеличении числа готовых к работе потоков планировщику пришлось бы делить процессорное время между ними, так как дополнительные потоки сами по себе не создают вычислительных ресурсов.

\section{Выводы}
Реализовано вычисление определителя по определению с распределением перестановок между потоками. Матрица читается из TXT, число исполнителей задаётся ключом запуска. Результаты собираются после завершения потоков. Проверяются ошибки выделений, потоков и арифметики. Трейс показывает создание двух рабочих потоков, ожидание и вывод ответа.

Для матрицы $13\times13$ все запуски дали определитель 7736. В переданных измерениях лучший результат получен при 12 исполнителях: 41.237 с вместо 244.304 с, ускорение 5.924. Основной выигрыш уже достигается на шести исполнителях, а увеличение до двенадцати даёт небольшой дополнительный эффект. Потоки ускоряют перебор, но факториальная сложность остаётся.

\section{Источники}
\begin{enumerate}
\item \href{https://man7.org/linux/man-pages/man3/pthread_create.3.html}{Linux man-pages: pthread\_create(3).}
\item \href{https://man7.org/linux/man-pages/man3/pthread_join.3.html}{Linux man-pages: pthread\_join(3).}
\item \href{https://man7.org/linux/man-pages/man2/clone.2.html}{Linux man-pages: clone(2), clone3.}
\item \href{https://man7.org/linux/man-pages/man2/futex.2.html}{Linux man-pages: futex(2).}
\item \href{https://man7.org/linux/man-pages/man2/clock_gettime.2.html}{Linux man-pages: clock\_gettime(2).}
\item \href{https://man7.org/linux/man-pages/man1/strace.1.html}{Руководство strace(1).}
\end{enumerate}

\appendixsection{А}{Исходный код программы и сборка}
\begin{listingblock}{Файл src/main.c}
\inputcode{../src/main.c}
\end{listingblock}
\begin{listingblock}{Файл src/matrix.h}
\inputcode{../src/matrix.h}
\end{listingblock}
\begin{listingblock}{Файл src/matrix.c}
\inputcode{../src/matrix.c}
\end{listingblock}
\begin{listingblock}{Файл src/matrix\_txt.h}
\inputcode{../src/matrix_txt.h}
\end{listingblock}
\begin{listingblock}{Файл src/matrix\_txt.c}
\inputcode{../src/matrix_txt.c}
\end{listingblock}
\begin{listingblock}{Файл src/permutation.h}
\inputcode{../src/permutation.h}
\end{listingblock}
\begin{listingblock}{Файл src/permutation.c}
\inputcode{../src/permutation.c}
\end{listingblock}
\begin{listingblock}{Файл src/determinant.h}
\inputcode{../src/determinant.h}
\end{listingblock}
\begin{listingblock}{Файл src/determinant.c}
\inputcode{../src/determinant.c}
\end{listingblock}
\begin{listingblock}{Файл src/CMakeLists.txt}
\inputcode{../src/CMakeLists.txt}
\end{listingblock}
\begin{listingblock}{Файл report/Dockerfile}
\inputcode{Dockerfile}
\end{listingblock}
\begin{listingblock}{Файл report/Makefile}
\inputcode{Makefile}
\end{listingblock}

\appendixsection{Б}{Полный вывод strace}
\begin{listingblock}{Трассировка программы с тремя исполнителями}
\inputterminal[numbers=left,numbersep=6pt,fontsize=\scriptsize]{inc/strace-3x3.txt}
\end{listingblock}

\appendixsection{В}{Исходные результаты измерений}
\begin{listingblock}{Ручные запуски для матрицы 13 на 13}
\inputterminal{inc/benchmarks.txt}
\end{listingblock}

\end{document}
